\documentclass[11pt]{article}
\usepackage[utf8]{inputenc}
\usepackage[T1]{fontenc}
\usepackage[italian]{babel}
\usepackage{geometry}
\geometry{a4paper, margin=1in}
\usepackage{xcolor}
\usepackage{listings}
\usepackage{hyperref}

\usepackage{tcolorbox}
\tcbuselibrary{breakable}
\begin{document}

\begin{tcolorbox}[colback=blue!5!white,colframe=blue!75!black,title=\textbf{DOMANDA}, breakable]
Esercizio su Web Socket
\begin{itemize}
    \item Realizzazione di una chat eseguita dentro un browser web
    \begin{itemize}
        \item Fare riferimento al materiale apposito condiviso con gli studenti
    \end{itemize}
\end{itemize}
\end{tcolorbox}

\begin{tcolorbox}[colback=green!5!white,colframe=green!75!black,title=\textbf{RISPOSTA}, breakable]
\textbf{Soluzione Esercizio 28:} Web Socket Chat

\textbf{Operazioni Preliminari}\\
Ho copiato l'intera cartella \texttt{WebSocket\_Chat} fornita nel materiale didattico all'interno della directory \texttt{es28 - Web Socket Chat}. Questo progetto non è scritto in C, bensì sfrutta Node.js e la libreria \texttt{Socket.IO} per fornire funzionalità Real-Time.

\textbf{Analisi dell'Architettura}\\
L'applicazione è composta da due parti fondamentali che comunicano tra loro:

\begin{enumerate}
    \item \textbf{Backend (Server in Node.js):}
    \begin{itemize}
        \item Il file \texttt{server.js} utilizza il framework \texttt{express} per avviare un server in ascolto sulla porta 4000.
        \item La prima cosa che fa il server è servire i file statici (l'interfaccia utente) presenti nella cartella \texttt{public}.
        \item Successivamente, aggancia a questo server il modulo \texttt{socket.io}. Da questo momento il server è in grado di accettare richieste di "Protocol Upgrade": quando un client lo richiede, la connessione \texttt{HTTP} standard viene promossa a connessione WebSocket (un tunnel \texttt{TCP} persistente e bidirezionale).
        \item Quando il server riceve un evento chiamato \texttt{message} da un qualsiasi client, usa l'istruzione \texttt{io.sockets.emit('UploadChat', data)} per prendere quel messaggio e "spargere la voce" (Broadcast) ritrasmettendolo istantaneamente a TUTTI i client collegati.
    \end{itemize}

    \item \textbf{Frontend (Client in Javascript):}
    \begin{itemize}
        \item All'interno di \texttt{public/chat.js}, il client apre una connessione WebSocket verso il server: \texttt{io.connect('http://localhost:4000')}.
        \item Quando l'utente preme il tasto per inviare, il browser "emette" l'evento \texttt{message} verso il server.
        \item Contestualmente, il browser si mette "in ascolto" di eventuali eventi \texttt{UploadChat} provenienti dal server. Appena ne riceve uno, il javascript manipola il DOM (aggiungendo codice HTML) per mostrare a schermo il nuovo messaggio ricevuto, aggiornando la chat in tempo reale senza dover ricaricare la pagina!
    \end{itemize}
\end{enumerate}

\textbf{Comandi per testare la Chat:}
\begin{lstlisting}[language=bash, basicstyle=\ttfamily\small]
# Posizionati nella cartella
cd "es28 - Web Socket Chat"

# Installa le dipendenze (se non sono gia' presenti nella cartella node_modules)
npm install

# Avvia il server
node server.js

# Ora apri il tuo browser preferito e apri DUE schede diverse puntandole a:
http://localhost:4000

# Prova a scrivere un messaggio in una scheda e vedrai che comparira' magicamente in tempo reale anche sull'altra!
\end{lstlisting}
\end{tcolorbox}

\end{document}
